\section[Rédaction Bourbaki \redac{239} : un faisceau mou qui n'est pas fin]{Rédaction Bourbaki \redac{239} : un faisceau mou qui n'est pas fin\protect\piste\protect\verifie{Redaction239}}\label{sec:bourbaki-239}

La rédaction \redac{239}, \q{Livre~VII. Topologie préhomologique. Chapitre~I. Faisceaux (état~2)}, des archives Bourbaki (Archives Henri Poincaré, \ahp{777}) est datée de juin 1956 par sa notice, qui la donne à Jean Dieudonné ; elle compte 102 pages. Elle définit les faisceaux mous et les faisceaux fins, démontre qu'un faisceau fin est mou, et le rédacteur écrit à la page~86 qu'il ne sait pas démontrer la réciproque, qu'il en doute, et qu'il n'a pas de contre-exemple. Sa transcription reste locale pour l'instant ; on n'en cite ici aucun passage.

\emph{Mou et fin.} Un faisceau $\mathcal F$ sur un espace $X$ associe à chaque ouvert $U$ un ensemble de sections $\mathcal F(U)$, avec des restrictions, de sorte que des sections qui se recollent sur un recouvrement proviennent d'une seule. Le faisceau est \emph{mou} (§~4, n\textsuperscript{o}\,3, p.~78) si une section définie au voisinage d'un fermé coïncide, près de ce fermé, avec une section globale : les sections se prolongent. Un faisceau de modules est \emph{fin} (§~4, n\textsuperscript{o}\,6, p.~85) s'il a assez d'endomorphismes pour découper : pour un fermé $F$ et un voisinage ouvert $V$ de $F$, il existe un endomorphisme du faisceau qui est l'identité près de $F$ et nul près du complémentaire de $V$. L'exemple type est le faisceau des fonctions continues réelles, où l'on multiplie par une fonction plateau ; la proposition~4 (p.~85--86) montre qu'un faisceau fin est mou, en découpant une section avec ces endomorphismes.

\emph{L'idée du contre-exemple.} Être fin demande des endomorphismes qui valent l'identité à un endroit et zéro à un autre, et un endomorphisme de groupes doit respecter la torsion. Le faisceau des applications continues à valeurs dans le cercle $\mathbf T$ a une section d'ordre~2, la constante $-1$. Son image par un endomorphisme est une application de carré~$1$, à valeurs dans $\{\pm 1\}$, donc constante sur un espace connexe : elle ne peut valoir $-1$ à un endroit et $1$ à un autre. Le même faisceau, sur $[0, 1]$, est pourtant mou : une application dans le cercle se prolonge, parce qu'on peut suivre son argument par petits morceaux.

\begin{encadre}
\begin{contrexemple}[réponse à la question de la rédaction \redac{239}, p.~86]
Sur $X = [0, 1]$, le faisceau $U \mapsto C(U, \mathbf T)$ des applications continues à valeurs dans le cercle, faisceau de groupes abéliens (c'est-à-dire de $\mathbf Z$-modules), est mou et n'est pas fin. La partie \q{pas fin} vaut sur tout espace connexe.
\end{contrexemple}
\end{encadre}

\begin{proof}
\emph{Pas fin.} Soit $(\varphi_U)$ une famille d'endomorphismes des groupes $C(U, \mathbf T)$, compatible aux restrictions, qui soit l'identité sur un voisinage ouvert d'un point $p$ et triviale sur un voisinage ouvert d'un point $q$. L'image $g = \varphi_X(-1)$ vérifie $g^2 = \varphi_X(1) = 1$ ; elle est donc à valeurs dans $\{\pm 1\}$, et constante puisque $X$ est connexe. Près de $p$, elle vaut $-1$ ; près de $q$, elle vaut $1$. C'est absurde. Pour $X = [0, 1]$, il suffit de prendre $F = \{0\}$ et $V = [0, 1[$.

\emph{Mou.} Soient $K$ un fermé de $[0, 1]$, $U$ un ouvert qui le contient et $f : [0, 1] \to \mathbf T$ une application continue sur $U$ (ses valeurs ailleurs sont quelconques). On découpe $[0, 1]$ en $M$ intervalles de longueur $1/M$, assez fins pour que tout intervalle qui rencontre le $1/M$-voisinage de $K$ soit contenu dans $U$ et que $f$ y varie de moins de $1$ ; on appelle ces intervalles les bons. Sur un bon intervalle $[x_i, x_{i+1}]$, le quotient $f(t)/f(x_i)$ reste dans le demi-plan droit, et son argument $\Delta_i(t)$ est une fonction continue ; on la prolonge à tout $[0, 1]$ en bloquant $t$ aux bornes de l'intervalle. Sur un mauvais intervalle, on fait tourner $\Delta_i$ linéairement de $0$ à l'argument de $f(x_{i+1})/f(x_i)$. L'application globale est $g = f(0)\, e^{i\sum_i \Delta_i}$. Le produit se télescope, et sur chaque bon intervalle il redonne $f$ ; comme les bons intervalles couvrent un voisinage de $K$, $g$ coïncide avec $f$ près de $K$ (figure~\ref{fig:bourbaki-239}). La construction n'utilise pas la théorie des revêtements.
\end{proof}

\begin{figure}[htbp]
\centering
\begin{tikzpicture}[every node/.style={font=\small}]
\fill[black!7] (0.77,-1.45) rectangle (3.63,1.95);
\fill[black!7] (6.27,-1.45) rectangle (8.58,1.95);
\node[black!60, above] at (2.2,1.95) {$U$};
\node[black!60, above] at (7.42,1.95) {$U$};
\draw[black!55] (0,-1.45) -- (11,-1.45);
\foreach \j in {0,...,20} { \draw[black!55] ({0.55*\j},-1.52) -- ({0.55*\j},-1.38); }
\draw[piste, line width=2.4pt] (1.65,-1.45) -- (2.75,-1.45);
\draw[piste, line width=2.4pt] (7.15,-1.45) -- (7.70,-1.45);
\node[piste, below] at (2.2,-1.52) {$K$};
\node[piste, below] at (7.42,-1.52) {$K$};
\node[below] at (0,-1.52) {$0$};
\node[below] at (11,-1.52) {$1$};
\draw[frise!55, line width=2.6pt, line join=round] (0.000,-0.330) -- (0.550,0.495);
\draw[frise!55, line width=2.6pt, line join=round] (0.550,0.495) -- (1.100,0.647);
\draw[frise!55, line width=2.6pt, line join=round] (1.100,0.647) -- (1.169,0.656) -- (1.238,0.664) -- (1.306,0.670) -- (1.375,0.675) -- (1.444,0.678) -- (1.513,0.680) -- (1.581,0.682) -- (1.650,0.682);
\draw[frise!55, line width=2.6pt, line join=round] (1.650,0.682) -- (1.719,0.681) -- (1.788,0.680) -- (1.856,0.678) -- (1.925,0.676) -- (1.994,0.673) -- (2.062,0.670) -- (2.131,0.666) -- (2.200,0.662);
\draw[frise!55, line width=2.6pt, line join=round] (2.200,0.662) -- (2.269,0.657) -- (2.338,0.653) -- (2.406,0.648) -- (2.475,0.643) -- (2.544,0.638) -- (2.612,0.632) -- (2.681,0.626) -- (2.750,0.620);
\draw[frise!55, line width=2.6pt, line join=round] (2.750,0.620) -- (2.819,0.613) -- (2.888,0.606) -- (2.956,0.598) -- (3.025,0.590) -- (3.094,0.580) -- (3.162,0.571) -- (3.231,0.560) -- (3.300,0.548);
\draw[frise!55, line width=2.6pt, line join=round] (3.300,0.548) -- (3.850,0.660);
\draw[frise!55, line width=2.6pt, line join=round] (3.850,0.660) -- (4.400,1.650);
\draw[frise!55, line width=2.6pt, line join=round] (4.400,1.650) -- (4.950,0.770);
\draw[frise!55, line width=2.6pt, line join=round] (4.950,0.770) -- (5.500,-0.550);
\draw[frise!55, line width=2.6pt, line join=round] (5.500,-0.550) -- (6.050,0.110);
\draw[frise!55, line width=2.6pt, line join=round] (6.050,0.110) -- (6.600,-1.077);
\draw[frise!55, line width=2.6pt, line join=round] (6.600,-1.077) -- (6.669,-1.080) -- (6.738,-1.078) -- (6.806,-1.073) -- (6.875,-1.064) -- (6.944,-1.050) -- (7.012,-1.033) -- (7.081,-1.011) -- (7.150,-0.986);
\draw[frise!55, line width=2.6pt, line join=round] (7.150,-0.986) -- (7.219,-0.957) -- (7.287,-0.924) -- (7.356,-0.888) -- (7.425,-0.848) -- (7.494,-0.805) -- (7.562,-0.760) -- (7.631,-0.711) -- (7.700,-0.660);
\draw[frise!55, line width=2.6pt, line join=round] (7.700,-0.660) -- (7.769,-0.607) -- (7.837,-0.551) -- (7.906,-0.494) -- (7.975,-0.435) -- (8.044,-0.375) -- (8.113,-0.314) -- (8.181,-0.253) -- (8.250,-0.191);
\draw[frise!55, line width=2.6pt, line join=round] (8.250,-0.191) -- (8.800,-0.880);
\draw[frise!55, line width=2.6pt, line join=round] (8.800,-0.880) -- (9.350,0.440);
\draw[frise!55, line width=2.6pt, line join=round] (9.350,0.440) -- (9.900,-1.045);
\draw[frise!55, line width=2.6pt, line join=round] (9.900,-1.045) -- (10.450,-0.220);
\draw[frise!55, line width=2.6pt, line join=round] (10.450,-0.220) -- (11.000,1.045);
\draw[black!45, thin] (0.000,-0.330) -- (0.275,-0.503) -- (0.550,0.495) -- (0.660,0.406) -- (0.770,0.577);
\draw[black!45, thin] (3.630,0.475) -- (3.740,-0.111) -- (3.850,0.660) -- (4.125,0.584) -- (4.400,1.650) -- (4.675,0.648) -- (4.950,0.770) -- (5.225,-1.015) -- (5.500,-0.550) -- (5.775,-1.171) -- (6.050,0.110) -- (6.160,1.277) -- (6.270,-1.011);
\draw[black!45, thin] (8.580,0.104) -- (8.690,-0.440) -- (8.800,-0.880) -- (9.075,-0.514) -- (9.350,0.440) -- (9.625,1.747) -- (9.900,-1.045) -- (10.175,0.187) -- (10.450,-0.220) -- (10.725,1.274) -- (11.000,1.045);
\draw[black, semithick] (0.770,0.577) -- (0.842,0.596) -- (0.913,0.613) -- (0.985,0.627) -- (1.056,0.640) -- (1.128,0.651) -- (1.199,0.660) -- (1.270,0.667) -- (1.342,0.672) -- (1.413,0.677) -- (1.485,0.680) -- (1.557,0.681) -- (1.628,0.682) -- (1.700,0.682) -- (1.771,0.681) -- (1.843,0.679) -- (1.914,0.676) -- (1.986,0.673) -- (2.057,0.670) -- (2.128,0.666) -- (2.200,0.662) -- (2.272,0.657) -- (2.343,0.653) -- (2.415,0.647) -- (2.486,0.642) -- (2.558,0.636) -- (2.629,0.631) -- (2.701,0.624) -- (2.772,0.618) -- (2.844,0.610) -- (2.915,0.603) -- (2.987,0.594) -- (3.058,0.585) -- (3.129,0.575) -- (3.201,0.565) -- (3.272,0.553) -- (3.344,0.540) -- (3.415,0.526) -- (3.487,0.510) -- (3.558,0.493) -- (3.630,0.475);
\draw[black, semithick] (6.270,-1.011) -- (6.328,-1.029) -- (6.385,-1.044) -- (6.443,-1.056) -- (6.501,-1.066) -- (6.559,-1.074) -- (6.616,-1.078) -- (6.674,-1.080) -- (6.732,-1.079) -- (6.790,-1.075) -- (6.847,-1.068) -- (6.905,-1.058) -- (6.963,-1.046) -- (7.021,-1.030) -- (7.079,-1.012) -- (7.136,-0.991) -- (7.194,-0.968) -- (7.252,-0.942) -- (7.309,-0.913) -- (7.367,-0.882) -- (7.425,-0.848) -- (7.483,-0.813) -- (7.540,-0.775) -- (7.598,-0.735) -- (7.656,-0.693) -- (7.714,-0.650) -- (7.772,-0.605) -- (7.829,-0.558) -- (7.887,-0.510) -- (7.945,-0.461) -- (8.003,-0.411) -- (8.060,-0.361) -- (8.118,-0.309) -- (8.176,-0.258) -- (8.233,-0.205) -- (8.291,-0.153) -- (8.349,-0.101) -- (8.407,-0.049) -- (8.465,0.003) -- (8.522,0.054) -- (8.580,0.104);
\fill[black!70] (0.000,-0.330) circle (1.5pt);
\fill[black!70] (0.550,0.495) circle (1.5pt);
\fill[black!70] (3.850,0.660) circle (1.5pt);
\fill[black!70] (4.400,1.650) circle (1.5pt);
\fill[black!70] (4.950,0.770) circle (1.5pt);
\fill[black!70] (5.500,-0.550) circle (1.5pt);
\fill[black!70] (6.050,0.110) circle (1.5pt);
\fill[black!70] (8.800,-0.880) circle (1.5pt);
\fill[black!70] (9.350,0.440) circle (1.5pt);
\fill[black!70] (9.900,-1.045) circle (1.5pt);
\fill[black!70] (10.450,-0.220) circle (1.5pt);
\fill[black!70] (11.000,1.045) circle (1.5pt);
\end{tikzpicture}
\caption{La preuve de \q{mou} sur un exemple, avec $M = 20$. En abscisse, $[0, 1]$ et ses points de subdivision ; en ordonnée, un argument, c'est-à-dire un relevé dans $\mathbf R$ des valeurs dans le cercle. Le fermé $K$ est en rouge, l'ouvert $U$ en gris. L'application $f$ est continue sur $U$ (trait noir) et quelconque ailleurs (trait gris) ; les points gris sont ses valeurs aux points de subdivision hors de $U$. L'application globale $g$ (trait bleu large) coïncide avec $f$ sur les bons intervalles, ceux qui rencontrent le $1/M$-voisinage de $K$, et tourne à vitesse constante sur les autres, d'une valeur de $f$ à la suivante par le plus court arc.}
\label{fig:bourbaki-239}
\end{figure}

\begin{sloppypar}
\noindent\emph{Vérification en Lean.} \q{Pas fin} est \texttt{not\_exists\_endomorphism}, sur tout espace préconnexe, avec la restriction \texttt{Opens.\allowbreak{}res}, et \texttt{not\_fine\_unitInterval} pour $[0, 1]$, $F = \{0\}$ et $V = [0, 1[$ ; \q{mou} est \texttt{exists\_eqOn\_nhds} et \texttt{exists\_eqOn\_nhds'}, dont l'application globale est \texttt{globalMap}, dans \href{\leanblobmain/Redaction239.lean}{\texttt{Redaction239.\allowbreak{}lean}}.
\end{sloppypar}

Ne sont pas démontrés ici : la construction du faisceau $U \mapsto C(U, \mathbf T)$ comme faisceau de mathlib, et les définitions de la rédaction (avec ses bases d'ouverts paracompacts et ses familles de fermés) ; les deux théorèmes vérifiés sont le contenu concret des deux propriétés pour ce faisceau. La proposition~4, \q{fin implique mou}, n'est pas formalisée.

\litterature{piste candidate. Le congrès avait noté la question comme ouverte (\emph{La Tribu}, n\textsuperscript{o}\,41). La question 423444 de MathOverflow (2022) demande un exemple de faisceau mou qui ne soit pas fin, et n'avait toujours pas de réponse le 1\textsuperscript{er} novembre 2026 \cite{mo423444}. Un contre-exemple imprimé n'a pas été trouvé : la recherche dans zbMATH Open (résumés et comptes rendus) et dans les titres et résumés d'arXiv n'a rien donné. R.~Godement, \emph{Topologie algébrique et théorie des faisceaux} (1958), II.3.7 et II.9, et G.~E.~Bredon, \emph{Sheaf Theory} (2\textsuperscript{e} éd., 1997), II.9 et ses exercices, n'ont pas pu être consultés légalement en ligne : la question reste à vérifier sur ces deux livres.}
